%%%%%
%
\section{S�lection dans un vecteur selon un second vecteur (�\textit{flag}�, \texttt{entier})}

\itemstyle{Op\'erateur~:} \textsf{select\_flag\_data}
\label{OP-select_flag_data}
\index{select\_flag\_data}

%%%%%OBJET%%%%
\itemstyle{Objet~:}
S�lection des donn�es d'un vecteur selon des
crit�res appliqu�s � un second vecteur de type \texttt{int} ("flag").


%%%%%ARGUMENTS%%%%
\bigskip
%%%%%ENTREES%%%%
\begin{tabarg}
\multicolumn{3}{l}{\titretab {Entr\'ees}}\\
\hline
$X$ & Vecteur, {\tt float} & S�rie de donn�es � filtrer (unit� quelconque).\\
flag & Vecteur, {\tt Int\_32} & vecteur sur lequel seront appliqu�s les crit�res (any)\\
$Crit1_{min}$ & Coeff., {\tt float} & valeur mini autoris�e pour la premi�re gamme valide\\
$Crit1_{max}$ & Coeff., {\tt float} & valeur maxi autoris�e pour la premi�re gamme valide\\
$Crit2_{min}$ & Coeff., {\tt float} & valeur mini autoris�e pour la deuxi�me gamme valide\\
$Crit2_{max}$ & Coeff., {\tt float} & valeur maxi autoris�e pour la deuxi�me gamme valide\\
\end{tabarg}

\medskip
%%%%%SORTIES%%%%
\begin{tabarg}
\multicolumn{3}{l}{\titretab {Sorties}}\\
\hline
$X_{sel}$ & Vecteur, {\tt float} & S�rie de donn�es filtr�es (m�me unit� que $X$)\\
\end{tabarg}


%%%%%FORMULE%%%%
\itemstyle{Formule}

%%select_flag_data
Si ((flag $>= Crit1_{min}$) et (flag $<= Crit1_{max}$) ) ou 
((flag $>= Crit2_{min}$) et (flag $<= Crit2_{max}$)) alors :

$X_{sel}[i] = X[i]$


Sinon, $X_{sel}[i] =$ mis\_val


%%%%%COMMENTAIRE%%%%
\itemstyle{Commentaire~:}
        Particularit� : les bornes de validit� sont pass�es sous
forme de \texttt{float} (m�me si elles doivent �tre appliqu�s � des \texttt{int})
pour pouvoir �tre NaN (pour signifier : pas de crit�re � appliquer).

  Dans le cas o� le vecteur-crit�re (�\textit{flag}�) est de type \texttt{float}, 
utiliser l'op�rateur \hyperref[OP-select_X_selon_Y]{\operfmt{select\_X\_selon\_Y}}.

%%%%%AUTEUR%%%%
\itemstyle{Auteur~:}
CNRM/GMEI/TRAMM

%%%%%REFERENCE%%%%
\itemstyle{R\'ef\'erences~:}

